\documentclass[a4paper, 11pt]{article}
\usepackage[
    top=0.75in,
    bottom=1in,
    left=0.75in,
    right=0.75in
]{geometry}
\usepackage{graphicx, tabularx, hyperref, amsmath, wrapfig}
\usepackage[usenames,dvipsnames]{color,xcolor}
\usepackage{listings}
\usepackage[siunitx]{circuitikz}
\usepackage{booktabs}
\usepackage{pgfplotstable, pgfplots}
\linespread{1.2}
\selectfont
\usepackage{xepersian}
% فونت گزارش‌کارهای قبلی «Dibaj» بوده است؛ برای جلوگیری از خطای کامپایل در سیستم‌هایی
% که این فونت نصب نیست، یک جانشینِ سازگار تعریف می‌کنیم.
\IfFontExistsTF{Dibaj}{
\settextfont[Scale=0.9]{Dibaj}
\setdigitfont[Scale=1]{Dibaj}
}{
\settextfont[Scale=1]{Amiri}
\setdigitfont[Scale=1]{Amiri}
}

\lstset{
    basicstyle=\ttfamily\small,
    keywordstyle=\color{blue},
    commentstyle=\color{green!50!black},
    stringstyle=\color{red!70!black},
    numbers=left,
    numberstyle=\tiny,
    stepnumber=1,
    breaklines=true,
    frame=single
}

\begin{document}


\begin{center}
\textbf{{\huge گزارش‌کار آزمایش شمارهٔ ۸:\\
بررسی اختلاف فاز در مدارهای AC با منحنی‌های لیساژو و تعیین $L$، $C$ و بسامد تشدید}}
\end{center}

\hrule
\noindent\begin{tabular}{rrl}
اعضای گروه: &  &  \\
 &  &  \\
\end{tabular}\hfill
\begin{tabular}{rr}
تاریخ انجام آزمایش: &  \\
دستیار آموزشی: &  \\
\end{tabular}\hfill
\begin{tabular}{rr}
گروه:  &   \\
زیرگروه: &  \lr{}  \\
\end{tabular}
\hrule

\section*{اهداف و خلاصهٔ شرح آزمایش}

هدف اصلی این آزمایش، آشنایی عملی با \textbf{روشِ منحنی‌های لیساژو} در اسیلوسکوپ (حالت \lr{XY}) و به‌کارگیری آن برای:
\begin{itemize}
    \item تعیین بسامد مجهول با مقایسه با یک بسامد مرجع و شمارش تعداد نقاط تماسِ منحنی با محورهای مختصات؛
    \item اندازه‌گیری اختلاف فاز بین دو سیگنال سینوسی هم‌فرکانس و استخراج $\phi$ از شکل بیضی لیساژو؛
    \item تعیین \textbf{ضریب خودالقایی} القاگر در مدار سری $RL$ با رسم نمودار $\tan\phi$ برحسب $f$؛
    \item تعیین \textbf{ظرفیت} خازن در مدار سری $RC$ با رسم نمودار $\tan\phi$ برحسب $1/f$؛
    \item بررسی مدار سری $RLC$ و مشاهدهٔ حالت \textbf{تشدید} از طریق تغییر علامت $\phi$ و تعیین بسامد تشدید $f_{\mathrm{res}}$.
\end{itemize}

در انجام آزمایش، مقاومت سریِ جعبه مقاومت روی مقدار ثابت
\[
R=300~\Omega
\]
تنظیم شد و در بخش ظرفیت، از خازن نامی
\[
C_{\mathrm{nom}}=10~\mu\mathrm{F}
\]
طبق دستورکار استفاده گردید.

\paragraph{خلاصهٔ نتایج}
\begin{itemize}
    \item با روش لیساژو و با مرجع $f_x=100~\mathrm{Hz}$ و نسبت تماس‌ها $\frac{N_x}{N_y}=\frac{2}{4}$، بسامد مجهول $f_y=50~\mathrm{Hz}$ به دست آمد.
    \item از برازش کمترین مربعاتِ \(\tan\phi\) برحسب \(f\) (مدار $RL$)، مقدار
    \[
    L=(0.69\pm 0.06)\ \mathrm{H}
    \]
    به دست آمد.
    \item از برازش کمترین مربعاتِ \(\tan\phi\) برحسب \(1/f\) (مدار $RC$)، مقدار
    \[
    C=(11.5\pm 0.2)\ \mu\mathrm{F}
    \]
    حاصل شد که با مقدار نامی \(10~\mu\mathrm{F}\) اختلاف نسبی حدود \(14.5\%\) دارد.
    \item در مدار $RLC$، با مشاهدهٔ تبدیل بیضی به خط راست (یعنی \(\phi=0\)) بسامد تشدید برابر
    \[
    f_{\mathrm{res}}=48.8~\mathrm{Hz}
    \]
    ثبت شد. با فرض \(C=10~\mu\mathrm{F}\)، از رابطهٔ تشدید \( \omega_{\mathrm{res}}^2=\frac{1}{LC}\) مقدار
    \[
    L_{\mathrm{res}}=(1.06\pm 0.02)\ \mathrm{H}
    \]
    به دست آمد که بزرگ‌تر از مقدار استخراج‌شده از مدار $RL$ است؛ این اختلاف در بخش «بحث و تحلیل» از دید منابع خطا و ناایده‌آلی عناصر تحلیل می‌شود.
\end{itemize}

\section{مبانی نظری و روابط مورد نیاز}

\subsection{منحنی لیساژو و استخراج اختلاف فاز}

دو سیگنال سینوسی هم‌فرکانس را در نظر می‌گیریم:
\[
x(t)=A\sin(\omega t),\qquad y(t)=B\sin(\omega t+\phi),
\]
که در حالت \lr{XY} اسیلوسکوپ، محور افقی و عمودی به‌ترتیب متناظر با \(x\) و \(y\) است. با حذف \(t\) می‌توان نشان داد که شکل کلی منحنی یک \textbf{بیضی} است:
\begin{equation}
\left(\frac{x}{A}\right)^2+\left(\frac{y}{B}\right)^2-2\left(\frac{x}{A}\right)\left(\frac{y}{B}\right)\cos\phi=\sin^2\phi.
\label{eq:lissajous_ellipse}
\end{equation}
برای استخراج \(\phi\) از اندازه‌گیری روی صفحهٔ اسیلوسکوپ، کافی است مقدار \(y\) در \(x=0\) را در نظر بگیریم. از \(x=0\Rightarrow \sin(\omega t)=0\) نتیجه می‌شود:
\[
y(t)=B\sin(\phi)\quad\Rightarrow\quad |y_0|=B|\sin\phi|.
\]
اگر \(B\) را از روی بیشینهٔ انحراف عمودی بیضی (نیم‌دامنهٔ عمودی) بخوانیم و \(y_0\) را از تقاطع بیضی با محور \(y\) به دست آوریم، داریم:
\begin{equation}
|\sin\phi|=\frac{|y_0|}{B}.
\label{eq:sinphi_from_screen}
\end{equation}
در این آزمایش (طبق دستورکار) تنظیم‌ها به گونه‌ای انجام شد که \(0<\phi<\frac{\pi}{2}\) باشد؛ در این بازه \(\sin\phi\) و \(\tan\phi\) مثبت‌اند و از رابطهٔ
\begin{equation}
\tan\phi=\frac{\sin\phi}{\sqrt{1-\sin^2\phi}}
\label{eq:tan_from_sin}
\end{equation}
می‌توان \(\tan\phi\) را استخراج کرد. (در مدارهای خازنی، علامت فیزیکی \(\phi\) منفی است؛ این نکته در بخش مربوط به مدار $RC$ اعمال می‌شود.)

\subsection{مدار سری $RL$}

برای مدار سری \(R\) و \(L\)، امپدانس القاگر \(Z_L=j\omega L\) و امپدانس کل:
\[
Z=R+j\omega L.
\]
اگر اختلاف فاز را به صورت \(\phi=\arg(Z)\) تعریف کنیم، خواهیم داشت:
\begin{equation}
\tan\phi=\frac{\omega L}{R}.
\label{eq:tanphi_RL}
\end{equation}
بنابراین نمودار \(\tan\phi\) برحسب \(f\) خطی است و شیب آن از رابطهٔ
\[
\tan\phi=\left(\frac{2\pi L}{R}\right)f
\]
به‌دست می‌آید.

\subsection{مدار سری $RC$}

برای مدار سری \(R\) و \(C\)، امپدانس خازن \(Z_C=-\dfrac{j}{\omega C}\) و امپدانس کل:
\[
Z=R-\frac{j}{\omega C}.
\]
در این حالت:
\begin{equation}
\tan\phi=-\frac{1}{\omega RC}.
\label{eq:tanphi_RC}
\end{equation}
پس اگر \(x=\dfrac{1}{f}\) را متغیر مستقل بگیریم:
\[
\tan\phi=-\left(\frac{1}{2\pi RC}\right)\frac{1}{f},
\]
و نمودار \(\tan\phi\) برحسب \(1/f\) خطی است.

\subsection{مدار سری $RLC$ و تشدید}

در مدار سری \(R\)–\(L\)–\(C\) داریم:
\[
Z=R+j\left(\omega L-\frac{1}{\omega C}\right),
\]
و بنابراین:
\begin{equation}
\tan\phi=\frac{\omega L-\frac{1}{\omega C}}{R}.
\label{eq:tanphi_RLC}
\end{equation}
شرط تشدید (کمینه شدن امپدانس و صفر شدن بخش موهومی) عبارت است از:
\begin{equation}
\omega_{\mathrm{res}}L-\frac{1}{\omega_{\mathrm{res}}C}=0
\quad\Rightarrow\quad
\omega_{\mathrm{res}}=\sqrt{\frac{1}{LC}}.
\label{eq:resonance_condition}
\end{equation}


\section{روش انجام آزمایش و تنظیمات دستگاه‌ها}

در تمام مراحل، اسیلوسکوپ در حالت \lr{XY} قرار داده شد (کانال افقی: \lr{X} و کانال عمودی: \lr{Y}) تا منحنی‌های لیساژو تشکیل شوند. برای پایدار شدن شکل، حساسیت عمودی/افقی (ولت بر تقسیم) و همچنین دامنهٔ سیگنال‌های ورودی به گونه‌ای تنظیم شدند که:
\begin{itemize}
    \item منحنی روی صفحه «بُرش نخورَد» و در محدودهٔ شبکه باقی بماند؛
    \item نسبت دامنه‌ها \(A\) و \(B\) به اندازهٔ کافی بزرگ باشد تا خواندن \(y_0\) و \(B\) با خطای کمتر انجام شود؛
    \item شکل بیضی کاملاً \textbf{ساکن} باشد (برای بخش تعیین بسامد مجهول، نسبت فرکانس‌ها باید گویا و نسبتاً ساده باشد).
\end{itemize}

در ادامه، برای هر بخش مدار متناظر بسته شد. برای شفافیت، شماتیک ایده‌آل مدارها در شکل‌های زیر آورده شده است.

\subsection{مدار مرجع برای بسامد مجهول (دو سیگنال مستقل)}

\begin{figure}[h!]
\centering
\begin{circuitikz}[scale=1]
\draw
(0,0) node[ocirc,label=left:{\lr{CH1 (X)}}] (X) {}
(0,-2) node[ocirc,label=left:{\lr{CH2 (Y)}}] (Y) {}
(3,0) node[draw,minimum width=2.2cm,minimum height=1cm,align=center] (gen1) {مولد 1\\$f_x$}
(3,-2) node[draw,minimum width=2.2cm,minimum height=1cm,align=center] (gen2) {مولد 2\\$f_y$}
;
\draw (X) -- (gen1.west);
\draw (Y) -- (gen2.west);
\end{circuitikz}
\caption{نمای شماتیک اتصال دو سیگنال مستقل به ورودی‌های \lr{X} و \lr{Y} برای تشکیل منحنی لیساژو و تعیین بسامد مجهول.}
\label{fig:setup_freq}
\end{figure}

\subsection{مدار سری $RL$}

\begin{figure}[h!]
\centering
\begin{circuitikz}[scale=1]
\draw
(0,0) to[sV,l={$V(t)$}] (0,3)
      to[R,l={$R=300\,\Omega$}] (3,3)
      to[L,l={$L$}] (6,3)
      -- (6,0) -- (0,0)
;
\draw (3,3) node[above] {$V_R(t)$};
\end{circuitikz}
\caption{مدار سری $RL$؛ اختلاف فاز بین \(V_R(t)\) (هم‌فاز با جریان) و ولتاژ منبع از طریق لیساژو اندازه‌گیری شد.}
\label{fig:RL_circuit}
\end{figure}

در این بخش، برای آن‌که اختلاف فاز در بازهٔ \(0<\phi<\pi/2\) قرار گیرد، فرکانس به گونه‌ای انتخاب شد که \(X_L=\omega L\) هم‌مرتبه با \(R\) باشد. سپس برای چند فرکانس مختلف، مقدار \(\sin\phi\) از نسبت \(y_0/B\) استخراج شد.

\subsection{مدار سری $RC$}

\begin{figure}[h!]
\centering
\begin{circuitikz}[scale=1]
\draw
(0,0) to[sV,l={$V(t)$}] (0,3)
      to[R,l={$R=300\,\Omega$}] (3,3)
      to[C,l={$C$}] (6,3)
      -- (6,0) -- (0,0)
;
\draw (3,3) node[above] {$V_R(t)$};
\end{circuitikz}
\caption{مدار سری $RC$؛ اختلاف فاز بین \(V_R(t)\) (هم‌فاز با جریان) و ولتاژ منبع اندازه‌گیری شد.}
\label{fig:RC_circuit}
\end{figure}

در این بخش نیز مشابه مدار $RL$، با تغییر فرکانس، اختلاف فاز اندازه‌گیری شد. چون در مدار خازنی جریان از ولتاژ جلوتر است، علامت \(\phi\) منفی است و باید پس از استخراج \(|\sin\phi|\) از لیساژو، علامت فیزیکی اعمال شود.

\subsection{مدار سری $RLC$}

\begin{figure}[h!]
\centering
\begin{circuitikz}[scale=1]
\draw
(0,0) to[sV,l={$V(t)$}] (0,3)
      to[R,l={$R=300\,\Omega$}] (3,3)
      to[L,l={$L$}] (5,3)
      to[C,l={$C=10\,\mu\mathrm{F}$}] (7,3)
      -- (7,0) -- (0,0)
;
\draw (3,3) node[above] {$V_R(t)$};
\end{circuitikz}
\caption{مدار سری $RLC$؛ با تغییر فرکانس، تغییر علامت اختلاف فاز و حالت تشدید مشاهده شد.}
\label{fig:RLC_circuit}
\end{figure}

برای مشاهدهٔ تشدید، فرکانس به‌تدریج تغییر داده شد تا در نقطه‌ای که \(X_L=X_C\) است (رابطهٔ \eqref{eq:resonance_condition}) اختلاف فاز صفر شود. در اسیلوسکوپ، این حالت به صورت تبدیل بیضی به \textbf{خط راست} (هم‌فاز شدن دو سیگنال) دیده می‌شود.

\section{بخش اول: تعیین بسامد مجهول با منحنی لیساژو}

در این بخش، دو سیگنال سینوسی با بسامدهای \(f_x\) و \(f_y\) به ورودی‌های افقی و عمودی اسیلوسکوپ داده شد. اگر نسبت بسامدها یک نسبت گویا باشد:
\[
\frac{f_x}{f_y}=\frac{p}{q},
\]
منحنی لیساژو \textbf{ساکن} می‌شود و تعداد نقاط تماس منحنی با محورهای \(x\) و \(y\) متناسب با همین نسبت است. در عمل، با شمارش تعداد نقاط تماس با محور \(x\) (به تعداد \(N_x\)) و محور \(y\) (به تعداد \(N_y\)) می‌توان نوشت:
\begin{equation}
\frac{f_y}{f_x}=\frac{N_x}{N_y}.
\label{eq:freq_from_lissajous}
\end{equation}

\begin{table}[h!]
\caption{اندازه‌گیری بسامد مجهول با روش لیساژو}
\label{tab:freq_unknown}
\centering
\begin{tabular}{ccc}
\toprule
بسامد مرجع \(f_x\) (Hz) & \(\dfrac{N_x}{N_y}\) & بسامد مجهول \(f_y\) (Hz)\\
\midrule
$100$ & $\dfrac{2}{4}$ & $50$ \\
\bottomrule
\end{tabular}
\end{table}

با جایگذاری داده‌های جدول \ref{tab:freq_unknown} در رابطهٔ \eqref{eq:freq_from_lissajous}:
\[
f_y=f_x\frac{N_x}{N_y}=100\times\frac{2}{4}=50~\mathrm{Hz}.
\]

\paragraph{یادداشت دربارهٔ عدم‌قطعیت}
در این روش، منبع اصلی خطا \textbf{شمارش نقاط تماس} است. چون \(N_x\) و \(N_y\) اعداد صحیح‌اند، در صورتی که شکل کاملاً پایدار نباشد (به دلیل ناپایداری فرکانس یا تنظیم نامناسب \lr{Trigger}) احتمال خطای \(\pm 1\) در شمارش وجود دارد. در دادهٔ حاضر نسبت \(\frac{2}{4}\) نسبت ساده‌ای است و خطای شمارش (اگر رخ دهد) به تغییر محسوس در نتیجه می‌انجامد؛ بنابراین در عمل، لازم است شکل با حوصله پایدار شود تا نسبت تماس‌ها بدون ابهام تعیین شود.

\section{بخش دوم: تعیین ضریب خودالقایی $L$ در مدار $RL$}

\subsection{داده‌های خام و استخراج \(\sin\phi\) و \(\tan\phi\)}

در این بخش، مدار سری \(R\)–\(L\) بسته شد و دو سیگنال متناظر با ولتاژ روی مقاومت و ولتاژ منبع (یا دو کمیت هم‌فرکانسِ مرجع و پاسخ) به ورودی‌های \lr{X} و \lr{Y} اسیلوسکوپ داده شد تا بیضی لیساژو تشکیل شود. طبق رابطهٔ \eqref{eq:sinphi_from_screen}:
\[
\sin\phi=\frac{y_0}{B},
\]
که در داده‌های خام، \(y_0/B\) به صورت «صورت/مخرج» گزارش شده است.

\begin{table}[h!]
\caption{داده‌های مدار $RL$ و محاسبهٔ \(\sin\phi\)، \(\phi\) و \(\tan\phi\)}
\label{tab:RL_raw}
\centering
\begin{tabular}{cccccc}
\toprule
$f$ (Hz) & $y_0$ & $B$ & $\sin\phi=y_0/B$ & $\phi$ (rad) & $\tan\phi$\\
\midrule
$30$  & $0.88$ & $1.92$ & $0.4583$ & $0.4769$ & $0.5157$ \\
$60$  & $1.48$ & $2.08$ & $0.7115$ & $0.7931$ & $1.0127$ \\
$90$  & $1.76$ & $2.13$ & $0.8263$ & $0.9731$ & $1.4670$ \\
$120$ & $1.88$ & $2.26$ & $0.8319$ & $0.9870$ & $1.4989$ \\
\bottomrule
\end{tabular}
\end{table}

محاسبهٔ ستون \(\phi\) با \(\phi=\arcsin(\sin\phi)\) و محاسبهٔ ستون \(\tan\phi\) با رابطهٔ \eqref{eq:tan_from_sin} انجام شد.

\subsection{برازش خطی \(\tan\phi\) برحسب \(f\) و استخراج \(L\)}

طبق رابطهٔ \eqref{eq:tanphi_RL}:
\[
\tan\phi=\left(\frac{2\pi L}{R}\right)f.
\]
از آن‌جا که در حد \(f\to 0\) انتظار داریم \(\tan\phi\to 0\)، مدل برازش را \textbf{خطیِ عبوری از مبدأ} در نظر می‌گیریم:
\begin{equation}
y = a f,\qquad y=\tan\phi.
\label{eq:fit_through_origin}
\end{equation}

\paragraph{فرمول کمترین مربعاتِ عبوری از مبدأ}
اگر \(n\) نقطه \((f_i,y_i)\) داشته باشیم، برآورد کمترین مربعات برای \(a\) برابر است با:
\begin{equation}
a=\frac{\sum_{i=1}^{n} f_i y_i}{\sum_{i=1}^{n} f_i^2}.
\label{eq:a_origin}
\end{equation}
باقیمانده‌ها \(r_i=y_i-af_i\) و
\[
\mathrm{SSE}=\sum r_i^2
\]
تعریف می‌شوند. در برازش یک‌پارامتری، برآورد واریانس خطا:
\[
s^2=\frac{\mathrm{SSE}}{n-1},
\]
و عدم‌قطعیت شیب از رابطهٔ استاندارد:
\begin{equation}
\sigma_a=\sqrt{\frac{s^2}{\sum f_i^2}}.
\label{eq:sigma_a_origin}
\end{equation}

\paragraph{محاسبات عددی}
برای داده‌های جدول \ref{tab:RL_raw} داریم:
\[
\sum f_i^2=30^2+60^2+90^2+120^2=27000,
\]
\[
\sum f_i y_i = 30(0.5157)+60(1.0127)+90(1.4670)+120(1.4989)=388.7.
\]
پس از \eqref{eq:a_origin}:
\[
a = \frac{388.7}{27000}=1.439\times 10^{-2}\ \mathrm{Hz^{-1}}.
\]
با محاسبهٔ باقیمانده‌ها، به‌دست آمد:
\[
\mathrm{SSE}=6.39\times 10^{-2},\qquad
s^2=\frac{6.39\times 10^{-2}}{3}=2.13\times 10^{-2},
\]
و بنابراین از \eqref{eq:sigma_a_origin}:
\[
\sigma_a=1.17\times 10^{-3}\ \mathrm{Hz^{-1}}.
\]

\paragraph{استخراج $L$ و انتشار عدم‌قطعیت}
از تطبیق \(a=\dfrac{2\pi L}{R}\) داریم:
\begin{equation}
L=\frac{aR}{2\pi}.
\label{eq:L_from_a}
\end{equation}
چون \(R\) در این آزمایش مقدار ثابتِ تنظیم‌شده است، عدم‌قطعیت غالب از \(\sigma_a\) می‌آید و:
\begin{equation}
\sigma_L=\left|\frac{\partial L}{\partial a}\right|\sigma_a=\frac{R}{2\pi}\sigma_a.
\label{eq:sigma_L}
\end{equation}

با \(R=300~\Omega\) نتیجه می‌شود:
\[
L=\frac{(1.439\times10^{-2})(300)}{2\pi}=0.686\ \mathrm{H},
\qquad
\sigma_L=\frac{300}{2\pi}(1.17\times10^{-3})=0.056\ \mathrm{H}.
\]
بنابراین:
\[
\boxed{L=(0.69\pm 0.06)\ \mathrm{H}.}
\]

\subsection{نمودار \(\tan\phi\) برحسب \(f\)}

\begin{figure}[h!]
\centering
\begin{tikzpicture}
\begin{axis}[
    width=11cm,
    height=6.2cm,
    xlabel={\lr{[$\mathrm{Hz}$]} فرکانس $f$},
    ylabel={$\tan\phi$},
    xmin=0, xmax=130,
    ymin=0, ymax=1.7,
    xmajorgrids=true,
    ymajorgrids=true,
    grid style=dashed,
    legend pos=north west,
]
\addplot[only marks, mark=square, mark size=2.7pt] coordinates {
    (30,0.5157)
    (60,1.0127)
    (90,1.4670)
    (120,1.4989)
};
\addplot[domain=0:130, samples=100] {0.01439*x};
\legend{داده‌ها, برازش $ \tan\phi = (1.439\times10^{-2})\, f$}
\end{axis}
\end{tikzpicture}
\caption{برازش خطی \(\tan\phi\) برحسب \(f\) برای مدار $RL$.}
\label{fig:tan_f_RL}
\end{figure}

\paragraph{تفسیر فیزیکی}
با افزایش فرکانس، راکتانس القایی \(X_L=\omega L\) افزایش می‌یابد و مدار از رفتار مقاومتی به رفتار القایی نزدیک می‌شود؛ بنابراین \(\phi\) و به تبع آن \(\tan\phi\) افزایشی است. نزدیک شدن \(\tan\phi\) به مقادیر بزرگ نشان‌دهندهٔ نزدیک شدن اختلاف فاز به \(\pi/2\) است.

\section{بخش سوم: تعیین ظرفیت $C$ در مدار $RC$}

\subsection{داده‌های خام و استخراج \(\tan\phi\)}

در مدار سری \(R\)–\(C\)، جریان از ولتاژ جلوتر است و اختلاف فاز \(\phi\) از دید تعریف \(\phi=\arg(Z)\) \textbf{منفی} است (رابطهٔ \eqref{eq:tanphi_RC}). از طرف دیگر، روش لیساژو تنها \(|\sin\phi|\) را می‌دهد. از این رو ابتدا \(|\sin\phi|\) و \(|\tan\phi|\) را مانند بخش قبل محاسبه کرده و سپس علامت منفی را اعمال می‌کنیم:
\[
\tan\phi = -|\tan\phi|.
\]

\begin{table}[h!]
\caption{داده‌های مدار $RC$ و محاسبهٔ \(\sin\phi\) و \(\tan\phi\)}
\label{tab:RC_raw}
\centering
\begin{tabular}{cccccc}
\toprule
$f$ (Hz) & $y_0$ & $B$ & $\sin\phi$ & $|\tan\phi|$ & $\tan\phi$\\
\midrule
$30$  & $1.72$ & $2.06$ & $0.8349$ & $1.5172$ & $-1.5172$ \\
$60$  & $1.28$ & $2.04$ & $0.6275$ & $0.8058$ & $-0.8058$ \\
$90$  & $0.96$ & $2.09$ & $0.4593$ & $0.5171$ & $-0.5171$ \\
$120$ & $0.76$ & $1.96$ & $0.3878$ & $0.4207$ & $-0.4207$ \\
\bottomrule
\end{tabular}
\end{table}

\subsection{برازش خطی \(\tan\phi\) برحسب \(1/f\) و استخراج \(C\)}

از رابطهٔ \eqref{eq:tanphi_RC} داریم:
\[
\tan\phi = -\left(\frac{1}{2\pi RC}\right)\frac{1}{f}.
\]
برای \(x=\frac{1}{f}\) (برحسب ثانیه) مدل خطی عبوری از مبدأ را به صورت
\[
y = m x,\qquad y=\tan\phi
\]
در نظر می‌گیریم. از مقایسه با رابطهٔ بالا:
\begin{equation}
m = -\frac{1}{2\pi RC}
\quad\Rightarrow\quad
C = -\frac{1}{2\pi R m}.
\label{eq:C_from_m}
\end{equation}

\paragraph{محاسبات عددی}
برای داده‌های جدول \ref{tab:RC_raw} (با \(x_i=1/f_i\))، از روابط \eqref{eq:a_origin} و \eqref{eq:sigma_a_origin} برای برازش عبوری از مبدأ استفاده می‌کنیم. نتیجهٔ برازش:
\[
m = -46.31\pm 0.82,
\]
که در آن \(m\) برحسب \(\mathrm{s^{-1}}\) (یعنی \(\tan\phi\) برحسب ثانیه) است.

پس از \eqref{eq:C_from_m} با \(R=300~\Omega\):
\[
C = -\frac{1}{2\pi(300)(-46.31)}=1.145\times 10^{-5}\ \mathrm{F}=11.45\ \mu\mathrm{F}.
\]
انتشار عدم‌قطعیت (با فرض ثابت بودن \(R\)) به صورت
\[
\frac{\sigma_C}{C}=\frac{\sigma_m}{|m|}
\]
خواهد بود، بنابراین:
\[
\sigma_C = 11.45\times\frac{0.82}{46.31}=0.20\ \mu\mathrm{F}.
\]
پس:
\[
\boxed{C=(11.5\pm 0.2)\ \mu\mathrm{F}.}
\]

\subsection{نمودار \(\tan\phi\) برحسب \(1/f\)}

\begin{figure}[h!]
\centering
\begin{tikzpicture}
\begin{axis}[
    width=11cm,
    height=6.2cm,
    xlabel={$1/f\ (\mathrm{s})$},
    ylabel={$\tan\phi$},
    xmin=0, xmax=0.04,
    ymin=-1.7, ymax=0.1,
    xmajorgrids=true,
    ymajorgrids=true,
    grid style=dashed,
    legend pos=north east,
]
\addplot[only marks, mark=square, mark size=2.7pt] coordinates {
    (0.0333333,-1.5172)
    (0.0166667,-0.8058)
    (0.0111111,-0.5171)
    (0.0083333,-0.4207)
};
\addplot[domain=0:0.04, samples=100] {-46.31*x};
\legend{داده‌ها, برازش $ \tan\phi = (-46.31)\,(1/f)$}
\end{axis}
\end{tikzpicture}
\caption{برازش خطی \(\tan\phi\) برحسب \(1/f\) برای مدار $RC$.}
\label{fig:tan_invf_RC}
\end{figure}

\subsection{مقایسه با مقدار نامی خازن و بحث کوتاه}

مقدار نامی خازن طبق دستورکار \(10~\mu\mathrm{F}\) بود، در حالی که مقدار به‌دست‌آمده \(11.5~\mu\mathrm{F}\) است. اختلاف نسبی:
\[
\frac{|C-C_{\mathrm{nom}}|}{C_{\mathrm{nom}}}\times 100
=
\frac{|11.5-10.0|}{10.0}\times 100
\simeq 14.5\%.
\]
این مقدار از مرتبهٔ تلورانس‌های رایج خازن‌های الکترولیتی (به‌ویژه در فرکانس پایین) و نیز اثر مقاومت سری معادل (\lr{ESR}) و نشتی خازن قابل توجیه است؛ ضمن آن‌که استخراج \(\phi\) از بیضی لیساژو در حالتی که \(\phi\) بزرگ است (نزدیک \(\pi/2\)) به‌طور محسوسی حساسیت خطا را افزایش می‌دهد.

\section{بخش چهارم: بررسی مدار $RLC$ و تعیین بسامد تشدید}

\subsection{داده‌های خام و تعیین علامت \(\tan\phi\)}

در مدار سری \(R\)–\(L\)–\(C\)، علامت \(\phi\) به این بستگی دارد که مدار در ناحیهٔ \textbf{خازنی} است یا \textbf{القایی}:
\[
\omega L-\frac{1}{\omega C}
\begin{cases}
<0 & \Rightarrow \phi<0 \ (\text{رفتار خازنی})\\
>0 & \Rightarrow \phi>0 \ (\text{رفتار القایی})
\end{cases}
\]
در داده‌ها، مقدار \(\sin\phi\) از بیضی لیساژو به‌صورت قدرمطلق گزارش می‌شود. برای اعمال علامت فیزیکی، از این واقعیت استفاده می‌کنیم که در فرکانس‌های کمتر از تشدید، مدار خازنی است و در فرکانس‌های بالاتر، مدار القایی می‌شود. نقطهٔ تشدید با مشاهدهٔ تبدیل بیضی به خط راست (یعنی \(\phi=0\)) ثبت شده است.

\begin{table}[h!]
\caption{داده‌های مدار $RLC$ و محاسبهٔ \(\tan\phi\) با علامت فیزیکی}
\label{tab:RLC_raw}
\centering
\begin{tabular}{cccccc}
\toprule
$f$ (Hz) & $y_0$ & $B$ & $|\sin\phi|$ & $|\tan\phi|$ & $\tan\phi$\\
\midrule
$20$  & $1.72$ & $2.00$ & $0.8600$ & $1.6853$ & $-1.6853$ \\
$35$  & $1.04$ & $1.94$ & $0.5361$ & $0.6350$ & $-0.6350$ \\
$48.8$& $0$    &  ---   & $0$      & $0$      & $0$ \\
$65$  & $0.88$ & $2.00$ & $0.4400$ & $0.4900$ & $+0.4900$ \\
$80$  & $1.32$ & $2.08$ & $0.6346$ & $0.8212$ & $+0.8212$ \\
\bottomrule
\end{tabular}
\end{table}

\subsection{تعیین بسامد تشدید و محاسبهٔ $L$ از شرط تشدید}

طبق جدول \ref{tab:RLC_raw}، در \(f=48.8~\mathrm{Hz}\) اختلاف فاز صفر شده است؛ بنابراین:
\[
f_{\mathrm{res}}=48.8~\mathrm{Hz}.
\]
با فرض این‌که مقدار خازن همان مقدار استفاده‌شده در مدار است و برابر مقدار نامی \(C=10~\mu\mathrm{F}\) (طبق دستورکار)، از رابطهٔ \eqref{eq:resonance_condition}:
\[
L=\frac{1}{\omega_{\mathrm{res}}^2 C}=\frac{1}{(2\pi f_{\mathrm{res}})^2 C}.
\]
پس:
\[
L_{\mathrm{res}}=\frac{1}{(2\pi\times 48.8)^2(10\times10^{-6})}=1.064\ \mathrm{H}.
\]

\paragraph{برآورد عدم‌قطعیت}
در این بخش، مقدار \(f_{\mathrm{res}}\) با تنظیم فرکانس مولد و مشاهدهٔ تبدیل بیضی به خط راست تعیین شده است. اگر عدم‌قطعیت تنظیم/خواندن فرکانس را به‌صورت محافظه‌کارانه \(\delta f_{\mathrm{res}}=\pm 0.5~\mathrm{Hz}\) در نظر بگیریم، از
\[
L\propto \frac{1}{f_{\mathrm{res}}^2}
\]
نتیجه می‌شود:
\[
\frac{\sigma_L}{L}=2\frac{\sigma_{f_{\mathrm{res}}}}{f_{\mathrm{res}}}
\quad\Rightarrow\quad
\sigma_{L_{\mathrm{res}}}=1.064\times 2\frac{0.5}{48.8}=0.022\ \mathrm{H}.
\]
بنابراین:
\[
\boxed{L_{\mathrm{res}}=(1.06\pm 0.02)\ \mathrm{H}.}
\]

\subsection{نمودار \(\tan\phi\) برحسب \(f\) و مشاهدهٔ عبور از صفر}

\begin{figure}[h!]
\centering
\begin{tikzpicture}
\begin{axis}[
    width=11cm,
    height=6.4cm,
    xlabel={\lr{[$\mathrm{Hz}$]} فرکانس $f$},
    ylabel={$\tan\phi$},
    xmin=15, xmax=85,
    ymin=-2.0, ymax=1.2,
    xmajorgrids=true,
    ymajorgrids=true,
    grid style=dashed,
    legend pos=south east,
]
\addplot[only marks, mark=square, mark size=2.7pt] coordinates {
    (20,-1.6853)
    (35,-0.6350)
    (48.8,0)
    (65,0.4900)
    (80,0.8212)
};
\addplot[domain=15:85, samples=2] {0}; % x-axis
\legend{داده‌ها}
\end{axis}
\end{tikzpicture}
\caption{تغییر \(\tan\phi\) برحسب \(f\) در مدار $RLC$ و مشاهدهٔ عبور از صفر در تشدید.}
\label{fig:tan_f_RLC}
\end{figure}


\section{کنترل سازگاری نتایج و مقایسهٔ پیش‌بینی‌ها}

یکی از راه‌های ارزیابی منطقی بودن نتایج، مقایسهٔ پیش‌بینی‌های متقاطع است. در این آزمایش سه کمیت \(L\)، \(C\) و \(f_{\mathrm{res}}\) به‌نوعی به هم مرتبط‌اند و باید از دید مرتبهٔ بزرگی و روند وابستگی‌ها با هم سازگار باشند.

\subsection{پیش‌بینی بسامد تشدید از \(L\) و \(C\) استخراج‌شده از فاز}

اگر از \(L\) به‌دست‌آمده از مدار $RL$ و \(C\) به‌دست‌آمده از مدار $RC$ استفاده کنیم، بسامد تشدید نظری از رابطهٔ \eqref{eq:resonance_condition}:
\[
f_0=\frac{1}{2\pi\sqrt{LC}}
\]
به دست می‌آید. با \(L=0.686~\mathrm{H}\) و \(C=1.145\times 10^{-5}~\mathrm{F}\):
\[
f_0=\frac{1}{2\pi\sqrt{(0.686)(1.145\times10^{-5})}}=56.8~\mathrm{Hz}.
\]
این مقدار از \(f_{\mathrm{res}}=48.8~\mathrm{Hz}\) بزرگ‌تر است.

\paragraph{انتشار خطا برای \(f_0\)}
چون \(f_0\propto (LC)^{-1/2}\)، داریم:
\[
\frac{\sigma_{f_0}}{f_0}=\frac{1}{2}\sqrt{\left(\frac{\sigma_L}{L}\right)^2+\left(\frac{\sigma_C}{C}\right)^2}.
\]
با \(\sigma_L=0.056\ \mathrm{H}\) و \(\sigma_C=0.20~\mu\mathrm{F}=2.0\times10^{-7}\ \mathrm{F}\):
\[
\frac{\sigma_L}{L}=0.0816,\qquad
\frac{\sigma_C}{C}=0.0175,
\]
پس:
\[
\frac{\sigma_{f_0}}{f_0}=\frac{1}{2}\sqrt{(0.0816)^2+(0.0175)^2}=0.0417
\quad\Rightarrow\quad
\sigma_{f_0}=56.8\times 0.0417=2.37~\mathrm{Hz}.
\]
بنابراین:
\[
f_0 = (56.8\pm 2.4)\ \mathrm{Hz}.
\]
اختلاف \(56.8\) با \(48.8\) حدود \(8~\mathrm{Hz}\) است که از مرتبهٔ چند برابرِ عدم‌قطعیت آماریِ بالا بزرگ‌تر است؛ پس احتمال وجود \textbf{خطای سیستماتیک} (در استخراج فاز یا ناایده‌آلی عناصر) تقویت می‌شود.

\subsection{پیش‌بینی \(L\) از روی \(f_{\mathrm{res}}\) و \(C_{\mathrm{nom}}\)}

در متن اصلی، با فرض \(C=10~\mu\mathrm{F}\) از دادهٔ تشدید، \(L_{\mathrm{res}}=1.06~\mathrm{H}\) به دست آمد. اگر به جای \(C_{\mathrm{nom}}\) از ظرفیت اندازه‌گیری‌شده \(C=11.5~\mu\mathrm{F}\) استفاده شود، مقدار \(L\) متناظر از تشدید:
\[
L=\frac{1}{(2\pi f_{\mathrm{res}})^2 C}
\]
خواهد بود. با \(f_{\mathrm{res}}=48.8~\mathrm{Hz}\) و \(C=11.5~\mu\mathrm{F}\):
\[
L=\frac{1}{(2\pi\times48.8)^2(11.5\times 10^{-6})}=0.925~\mathrm{H}.
\]
این مقدار همچنان از \(0.686~\mathrm{H}\) بزرگ‌تر است، اما اختلاف کاهش می‌یابد. لذا بخشی از اختلاف میان دو روش می‌تواند صرفاً ناشی از تفاوت ظرفیت واقعی با مقدار نامی باشد.

\subsection{کنترل ابعادی روابط}

برای جلوگیری از خطای مفهومی در روابط، کنترل واحدها انجام می‌دهیم:
\begin{itemize}
    \item در مدار $RL$: \(\tan\phi=\omega L/R\) بی‌بُعد است چون \([\omega]=\mathrm{s^{-1}}\)، \([L]=\mathrm{H}=\Omega\cdot \mathrm{s}\) و \([R]=\Omega\).
    \item در مدار $RC$: \(\tan\phi=-1/(\omega RC)\) بی‌بُعد است چون \([RC]=\Omega\cdot\mathrm{F}=\mathrm{s}\).
    \item در تشدید: \(\omega_{\mathrm{res}}=\sqrt{1/(LC)}\) چون \([LC]=(\Omega\cdot\mathrm{s})\cdot(\mathrm{s}/\Omega)=\mathrm{s^2}\).
\end{itemize}
این کنترل نشان می‌دهد روابط استفاده‌شده از نظر ابعادی سازگارند و اختلاف‌های عددی بیشتر به خطاهای تجربی و ناایده‌آلی‌ها برمی‌گردد.

\section{بحث و تحلیل منابع اختلاف و خطا}

\subsection{دقت استخراج \(\phi\) از منحنی لیساژو}

رابطهٔ \eqref{eq:sinphi_from_screen} نشان می‌دهد که \(\sin\phi\) از نسبت دو کمیت خوانده‌شده از روی صفحه به دست می‌آید. در عمل:
\begin{itemize}
    \item \textbf{تفکیک‌پذیری شبکهٔ اسیلوسکوپ} و ضخامت خطِ منحنی باعث می‌شود خواندن \(y_0\) و \(B\) با عدم‌قطعیت همراه باشد.
    \item وقتی \(\phi\) بزرگ باشد، \(\sin\phi\) به 1 نزدیک می‌شود و در نتیجه طبق \eqref{eq:tan_from_sin}، \(\tan\phi\) نسبت به تغییرات کوچک \(\sin\phi\) بسیار حساس می‌شود (چون مخرج \(\sqrt{1-\sin^2\phi}\) کوچک می‌شود). این اثر در داده‌های مدار $RL$ در فرکانس‌های بالاتر به‌وضوح دیده می‌شود.
\end{itemize}

\subsection{ناایده‌آلی عناصر: مقاومت سیم‌پیچ و \lr{ESR} خازن}

\paragraph{القاگر}
در استخراج رابطهٔ \eqref{eq:tanphi_RL} فرض شد که القاگر فقط راکتانس \(j\omega L\) دارد. اما یک سیم‌پیچ واقعی دارای مقاومت اهمی \(R_L\) نیز هست. در مدار سری واقعی:
\[
Z=(R+R_L)+j\omega L
\quad\Rightarrow\quad
\tan\phi=\frac{\omega L}{R+R_L}.
\]
اگر \(R_L\) قابل توجه باشد، استفاده از \(R\) به‌جای \(R+R_L\) سبب می‌شود \(L\) \textbf{کم‌تر از مقدار واقعی} برآورد شود؛ این نکته می‌تواند یکی از دلایل آن باشد که \(L\) به‌دست‌آمده از مدار $RL$ کوچک‌تر از \(L_{\mathrm{res}}\) محاسبه‌شده از تشدید است.

\paragraph{خازن}
در مورد خازن‌های الکترولیتی، \lr{ESR} و نشتیِ دی‌الکتریک می‌تواند موجب شود رفتار عنصر کاملاً \( -j/(\omega C)\) نباشد؛ این موضوع هم روی اختلاف فاز و هم روی شیب نمودار \(\tan\phi\) برحسب \(1/f\) اثر سیستماتیک می‌گذارد.

\subsection{چرا اختلاف $L$ از دو روش طبیعی است؟}

در این آزمایش دو راه مستقل برای برآورد \(L\) داشتیم:
\begin{itemize}
    \item روش اول: از رابطهٔ \(\tan\phi=\omega L/R\) در مدار $RL$.
    \item روش دوم: از شرط تشدید در مدار $RLC$ و فرکانس عبور فاز از صفر.
\end{itemize}
روش اول به‌طور مستقیم به \textbf{خواندن دقیق بیضی} حساس است، و علاوه بر آن، نسبت به مقاومت داخلی سیم‌پیچ نیز حساسیت دارد. روش دوم، اگر فرکانس تشدید با دقت خوبی قابل تنظیم/اندازه‌گیری باشد، از نظر مفهومی پایدارتر است؛ اما به مقدار واقعی خازن (و تغییر آن با فرکانس) نیز وابسته است. لذا اختلاف مشاهده‌شده به‌خودی‌خود غیرمنطقی نیست، ولی باید در گزارش با همین منطق توجیه و منابع خطا مشخص شوند.


\section*{پرسش‌ها و نکات تکمیلی}

\paragraph{۱) چرا در مدار القاگر، ولتاژ از جریان عقب/جلو می‌افتد؟}
برای القاگر ایده‌آل داریم \(V_L=L\,\dfrac{dI}{dt}\). اگر جریان \(I(t)=I_m\sin(\omega t)\) باشد،
\[
V_L(t)=L\omega I_m\cos(\omega t)=L\omega I_m\sin\!\left(\omega t+\frac{\pi}{2}\right),
\]
پس ولتاژِ القاگر \(\frac{\pi}{2}\) \textbf{از جریان جلوتر} است. در مدار سری \(R\)–\(L\)، حضور مقاومت باعث می‌شود اختلاف فاز کل بین \(0\) و \(\pi/2\) قرار گیرد و مقدار آن با \(\tan\phi=\omega L/R\) کنترل می‌شود.

\paragraph{۲) چرا در مدار خازن، جریان از ولتاژ جلوتر است؟}
برای خازن ایده‌آل داریم \(I=C\,\dfrac{dV}{dt}\). اگر ولتاژ \(V(t)=V_m\sin(\omega t)\) باشد،
\[
I(t)=C\omega V_m\cos(\omega t)=C\omega V_m\sin\!\left(\omega t+\frac{\pi}{2}\right),
\]
یعنی \textbf{جریان \(\frac{\pi}{2}\) از ولتاژ جلوتر} است. در مدار سری \(R\)–\(C\)، وجود مقاومت موجب می‌شود اختلاف فاز کل بین \(0\) و \(-\pi/2\) قرار گیرد و \(\tan\phi=-1/(\omega RC)\) برقرار باشد.

\paragraph{۳) در تشدید سری چه رخ می‌دهد و چرا \(\phi=0\) می‌شود؟}
در مدار سری \(R\)–\(L\)–\(C\)، بخش موهومی امپدانس برابر \(\omega L-1/(\omega C)\) است. در فرکانس تشدید:
\[
\omega_{\mathrm{res}}L=\frac{1}{\omega_{\mathrm{res}}C}\quad\Rightarrow\quad \Im(Z)=0,
\]
بنابراین امپدانس کل صرفاً مقاومتی است و جریان با ولتاژ \textbf{هم‌فاز} می‌شود، یعنی \(\phi=0\). از دید انرژی نیز در این حالت، انرژی بین میدان مغناطیسی القاگر و میدان الکتریکی خازن به تناوب ردوبدل می‌شود و تلفات صرفاً از مقاومت (و مؤلفه‌های اتلافی) ناشی می‌گردد.

\paragraph{۴) مهم‌ترین خطاهای آزمایش و علت آن‌ها}
\begin{itemize}
    \item \textbf{خطای خواندن \(y_0\) و \(B\)} از روی شبکهٔ اسیلوسکوپ (خطای دید و ضخامت منحنی).
    \item \textbf{حساسیت بالا در \(\phi\)‌های بزرگ}: وقتی \(\sin\phi\to 1\)، تبدیل به \(\tan\phi\) بسیار حساس می‌شود (ضمیمهٔ B).
    \item \textbf{مقاومت داخلی سیم‌پیچ} و \textbf{\lr{ESR} خازن} که روابط ایده‌آل فاز را تغییر می‌دهد.
    \item \textbf{پراکندگی پارامترها با فرکانس}: به‌ویژه در خازن‌های الکترولیتی، ظرفیت مؤثر و \lr{ESR} تابع فرکانس و دما است.
    \item \textbf{اتصالات و سیم‌کشی}: القای متقابل، مقاومت تماس‌ها و ظرفیت‌های پراکنده می‌تواند در فرکانس‌های پایین نیز اثر قابل‌ملاحظه داشته باشد.
\end{itemize}


\section*{جمع‌بندی}

با استفاده از منحنی‌های لیساژو در اسیلوسکوپ، هم بسامد مجهول و هم اختلاف فاز مدارهای AC تعیین شد. سپس با روابط امپدانس در مدارهای سری:
\begin{itemize}
    \item ضریب خودالقایی القاگر از برازش \(\tan\phi\) برحسب \(f\) در مدار $RL$ به صورت \(L=(0.69\pm 0.06)\,\mathrm{H}\) به دست آمد؛
    \item ظرفیت خازن از برازش \(\tan\phi\) برحسب \(1/f\) در مدار $RC$ به صورت \(C=(11.5\pm 0.2)\,\mu\mathrm{F}\) محاسبه شد؛
    \item در مدار $RLC$ بسامد تشدید \(f_{\mathrm{res}}=48.8~\mathrm{Hz}\) مشاهده شد و بر این اساس \(L_{\mathrm{res}}=(1.06\pm 0.02)\,\mathrm{H}\) به دست آمد.
\end{itemize}
تفاوت دو مقدار \(L\) با توجه به ناایده‌آلی عناصر (به‌ویژه مقاومت داخلی سیم‌پیچ و عدم‌قطعیت استخراج \(\phi\) از بیضی لیساژو) قابل توجیه است.


\appendix
\section{ضمیمهٔ A: محاسبات عددیِ باقیمانده‌ها و عدم‌قطعیت شیب}

در این ضمیمه، برای جلوگیری از «نوشتن عدد بدون استدلال»، محاسبات \lr{SSE} و باقیمانده‌ها به صورت صریح ارائه می‌شود.

\subsection{مدار $RL$}

برازش عبوری از مبدأ به صورت \(y=af\) با \(a=1.439\times 10^{-2}\) انجام شد. جدول \ref{tab:RL_resid} مقادیر \(y_i\)، \(af_i\) و باقیمانده‌ها را نشان می‌دهد.

\begin{table}[h!]
\caption{محاسبهٔ باقیمانده‌ها در مدار $RL$}
\label{tab:RL_resid}
\centering
\begin{tabular}{cccc}
\toprule
$f_i$ (Hz) & $y_i=\tan\phi$ & $af_i$ & $r_i=y_i-af_i$\\
\midrule
$30$  & $0.5157$ & $0.4317$ & $0.0840$\\
$60$  & $1.0127$ & $0.8634$ & $0.1493$\\
$90$  & $1.4670$ & $1.2951$ & $0.1719$\\
$120$ & $1.4989$ & $1.7268$ & $-0.2279$\\
\bottomrule
\end{tabular}
\end{table}

پس:
\[
\mathrm{SSE}=\sum r_i^2
=0.0840^2+0.1493^2+0.1719^2+(-0.2279)^2
=6.39\times 10^{-2}.
\]
و همان‌طور که در متن آمده بود:
\[
s^2=\frac{\mathrm{SSE}}{n-1}=\frac{6.39\times 10^{-2}}{3}=2.13\times 10^{-2}.
\]
اکنون از \eqref{eq:sigma_a_origin}:
\[
\sigma_a=\sqrt{\frac{s^2}{\sum f_i^2}}
=\sqrt{\frac{2.13\times 10^{-2}}{27000}}
=1.17\times 10^{-3}.
\]

\subsection{مدار $RC$}

برای مدار $RC$ برازش عبوری از مبدأ به صورت \(y=mx\) با \(x=1/f\) انجام شد و \(m=-46.31\) به دست آمد. جدول \ref{tab:RC_resid} باقیمانده‌ها را نشان می‌دهد.

\begin{table}[h!]
\caption{محاسبهٔ باقیمانده‌ها در مدار $RC$}
\label{tab:RC_resid}
\centering
\begin{tabular}{ccccc}
\toprule
$f_i$ (Hz) & $x_i=1/f_i$ (s) & $y_i=\tan\phi$ & $mx_i$ & $r_i$\\
\midrule
$30$  & $0.03333$ & $-1.5172$ & $-1.5437$ & $0.0265$\\
$60$  & $0.01667$ & $-0.8058$ & $-0.7718$ & $-0.0340$\\
$90$  & $0.01111$ & $-0.5171$ & $-0.5146$ & $-0.0025$\\
$120$ & $0.00833$ & $-0.4207$ & $-0.3859$ & $-0.0348$\\
\bottomrule
\end{tabular}
\end{table}

بنابراین:
\[
\mathrm{SSE}=0.0265^2+(-0.0340)^2+(-0.0025)^2+(-0.0348)^2
=3.07\times 10^{-3},
\]
و
\[
s^2=\frac{\mathrm{SSE}}{n-1}=\frac{3.07\times10^{-3}}{3}=1.02\times10^{-3}.
\]
همچنین:
\[
\sum x_i^2=(0.03333)^2+(0.01667)^2+(0.01111)^2+(0.00833)^2
=1.584\times 10^{-3}.
\]
پس:
\[
\sigma_m=\sqrt{\frac{s^2}{\sum x_i^2}}
=\sqrt{\frac{1.02\times10^{-3}}{1.584\times10^{-3}}}=0.82.
\]

\section{ضمیمهٔ B: انتشار خطا از خواندن \(y_0\) و \(B\)}

در صورت نیاز به یک برآورد سیستماتیک از خطای اندازه‌گیری \(\sin\phi\) ناشی از خواندن مقادیر روی صفحهٔ اسیلوسکوپ، می‌توان چنین عمل کرد. اگر:
\[
s\equiv \sin\phi=\frac{y_0}{B},
\]
و عدم‌قطعیت خواندن \(y_0\) و \(B\) به‌ترتیب \(\sigma_{y_0}\) و \(\sigma_B\) باشد، با فرض مستقل بودن خطاها:
\[
\sigma_s^2=\left(\frac{\partial s}{\partial y_0}\right)^2\sigma_{y_0}^2+\left(\frac{\partial s}{\partial B}\right)^2\sigma_B^2
=
\left(\frac{1}{B}\right)^2\sigma_{y_0}^2+\left(\frac{y_0}{B^2}\right)^2\sigma_B^2.
\]
اکنون چون \(\tan\phi=\dfrac{s}{\sqrt{1-s^2}}\)، مشتق \(\tan\phi\) نسبت به \(s\) برابر است با:
\[
\frac{d}{ds}\left(\frac{s}{\sqrt{1-s^2}}\right)=\frac{1}{(1-s^2)^{3/2}}.
\]
پس:
\[
\sigma_{\tan\phi}=\frac{\sigma_s}{(1-s^2)^{3/2}}.
\]
این رابطه نشان می‌دهد هرچه \(s\) به 1 نزدیک‌تر شود (یعنی \(\phi\to \pi/2\))، \(\sigma_{\tan\phi}\) به‌شدت بزرگ می‌شود. این دقیقاً همان چراییِ حساسیت بالای داده‌های فاز در اختلاف‌فازهای بزرگ است و توضیح می‌دهد چرا خطاهای خواندنِ ساده روی شبکهٔ اسیلوسکوپ می‌تواند در نهایت به اختلاف محسوس در \(L\) و \(C\) منجر شود.

\section{ضمیمهٔ C: استخراج علامت اختلاف فاز از شکل لیساژو}

روش \eqref{eq:sinphi_from_screen} فقط \(|\sin\phi|\) را می‌دهد. برای تعیین علامت \(\phi\) باید دانست کدام سیگنال از دیگری جلوتر است. در عمل دو راه رایج وجود دارد:
\begin{itemize}
    \item استفاده از حالت \lr{Time Domain} (نمایش برحسب زمان) و مقایسهٔ جابه‌جایی زمانی دو موج؛
    \item یا توجه به جهت گردش منحنی لیساژو در حالت \lr{XY} (با استفاده از حالت \lr{Persist} یا مشاهدهٔ حرکت نقطهٔ روشن) که نشان می‌دهد کدام کانال تقدم فاز دارد.
\end{itemize}
در این گزارش، چون داده‌های خام فقط نسبت‌های \(y_0/B\) را شامل می‌شد، علامت فیزیکی در مدار $RC$ و $RLC$ از دانش کیفیِ رفتار مدار (خازنی/القایی) تعیین شد.



\section{ضمیمهٔ D: استخراج معادلهٔ بیضی لیساژو}

از
\[
x=A\sin(\omega t),\qquad y=B\sin(\omega t+\phi)=B(\sin\omega t\cos\phi+\cos\omega t\sin\phi)
\]
داریم:
\[
\sin(\omega t)=\frac{x}{A},\qquad
\cos(\omega t)=\pm\sqrt{1-\sin^2(\omega t)}=\pm\sqrt{1-\left(\frac{x}{A}\right)^2}.
\]
همچنین:
\[
\frac{y}{B}=\frac{x}{A}\cos\phi+\cos(\omega t)\sin\phi.
\]
پس:
\[
\frac{y}{B}-\frac{x}{A}\cos\phi=\cos(\omega t)\sin\phi.
\]
با توان دوم گرفتن:
\[
\left(\frac{y}{B}-\frac{x}{A}\cos\phi\right)^2=\left(1-\left(\frac{x}{A}\right)^2\right)\sin^2\phi.
\]
گسترش و مرتب‌سازی دقیقاً به رابطهٔ \eqref{eq:lissajous_ellipse} می‌انجامد:
\[
\left(\frac{x}{A}\right)^2+\left(\frac{y}{B}\right)^2-2\left(\frac{x}{A}\right)\left(\frac{y}{B}\right)\cos\phi=\sin^2\phi.
\]

\section{ضمیمهٔ E: برازش با عرض از مبدأ و دلیل ترجیح مدل عبوری از مبدأ}

اگر به جای مدل عبوری از مبدأ، مدل کلی \(y=af+b\) را برای داده‌های مدار $RL$ در نظر بگیریم، از کمترین مربعات معمولی مقدارهای تقریبی
\[
a\simeq 1.135\times 10^{-2},\qquad b\simeq 0.327
\]
به دست می‌آید (اعداد تقریبی‌اند). اما از نظر فیزیکی:
\[
\lim_{f\to 0}\tan\phi=\lim_{f\to 0}\frac{\omega L}{R}=0.
\]
بنابراین انتظار داریم \(b=0\) باشد و هر عرض از مبدأ غیرصفر، بیشتر بازتاب محدودیت داده‌ها، خطای خواندن و ناایده‌آلی مدار است تا یک ویژگی فیزیکی واقعی. به همین دلیل، در این گزارش مدل عبوری از مبدأ به‌عنوان مدل اصلی به کار گرفته شد تا پارامتر \(L\) مستقیماً از شیب استخراج شود و تفسیر فیزیکی آن بدون ابهام باقی بماند.


\section{ضمیمهٔ F: نسبت فرکانس‌ها و تعداد تماس‌ها در منحنی لیساژو}

در بخش تعیین بسامد مجهول، از قاعدهٔ شمارش تماس‌ها استفاده شد. در این‌جا به‌صورت خلاصه توضیح می‌دهیم چرا رابطهٔ \eqref{eq:freq_from_lissajous} معقول است.

اگر:
\[
x(t)=A\sin(2\pi f_x t),\qquad y(t)=B\sin(2\pi f_y t+\phi),
\]
و نسبت فرکانس‌ها \(f_x/f_y=p/q\) باشد، در بازهٔ زمانی
\[
T=q\left(\frac{1}{f_x}\right)=p\left(\frac{1}{f_y}\right)
\]
هر دو سیگنال یک تعداد صحیحِ دوره کامل می‌زنند و منحنی بسته می‌شود. تعداد «لوب»‌ها و تماس‌ها با محورهای مختصات، از تعداد دفعاتی ناشی می‌شود که یکی از دو سیگنال به مقدار صفر می‌رسد در حالی که دیگری هنوز تغییر می‌کند. به بیان دیگر:
\begin{itemize}
    \item تماس با محور \(x\) زمانی رخ می‌دهد که \(y(t)=0\Rightarrow \sin(2\pi f_y t+\phi)=0\). در بازهٔ بسته شدن منحنی، این شرط \(2q\) بار برقرار می‌شود (با در نظر گرفتن صفرهای تابع سینوس در هر دوره).
    \item تماس با محور \(y\) زمانی رخ می‌دهد که \(x(t)=0\Rightarrow \sin(2\pi f_x t)=0\). در همان بازه، این شرط \(2p\) بار برقرار می‌شود.
\end{itemize}
در نتیجه نسبت تعداد تماس‌ها (با حذف ضریب 2) متناسب با \(q/p\) یا \(p/q\) می‌شود. این‌که در عمل از کدام نسبت استفاده می‌کنیم به این بستگی دارد که \(N_x\) را «تماس با کدام محور» تعریف کرده باشیم و نیز کدام سیگنال به کانال \(X\) متصل شده است. در این گزارش، تعریفِ \(N_x\) و \(N_y\) همان است که در دستورکار آمده و نتیجهٔ نهایی \(f_y=f_x(N_x/N_y)\) با داده‌های جدول \ref{tab:freq_unknown} سازگار است.

\paragraph{نکتهٔ عملی}
اگر نسبت \(p/q\) به اعداد بزرگ منجر شود، شکل لیساژو بسیار پیچیده می‌شود و شمارش تماس‌ها خطاپذیر است. بنابراین معمولاً بهتر است بسامد مرجع طوری تنظیم شود که نسبت فرکانس‌ها \textbf{ساده} باشد (مانند \(1/2\)، \(2/3\)، \(3/4\) و ...).

\paragraph{اثر اختلاف فاز}
وجود \(\phi\) شکل لیساژو را می‌چرخاند یا جابه‌جا می‌کند، اما \textbf{تعداد تماس‌ها} را (برای منحنی ساکن) تغییر نمی‌دهد؛ لذا روش تعیین بسامد مجهول نسبت به \(\phi\) حساسیت مستقیم ندارد. حساسیت اصلی همان پایدار بودن شکل و صحیح بودن شمارش است.



\section{ضمیمهٔ G: نمونهٔ محاسبهٔ گام‌به‌گام از دادهٔ خام تا \(\tan\phi\)}

برای روشن بودن روند محاسبات، یک نمونه از داده‌های مدار $RL$ (فرکانس \(60~\mathrm{Hz}\)) را کاملاً مرحله‌به‌مرحله انجام می‌دهیم.

طبق جدول \ref{tab:RL_raw} نسبت خوانده‌شده از اسیلوسکوپ:
\[
\frac{y_0}{B}=\frac{1.48}{2.08}.
\]
پس:
\[
\sin\phi=\frac{1.48}{2.08}=0.711538\ldots
\]
اکنون:
\[
\phi=\arcsin(0.711538\ldots)=0.7931\ \mathrm{rad}.
\]
برای محاسبهٔ \(\tan\phi\) دو راه معادل وجود دارد:
\begin{itemize}
    \item راه مستقیم: \(\tan\phi=\tan(0.7931)=1.0127\).
    \item راه مبتنی بر \(\sin\phi\): از \eqref{eq:tan_from_sin} داریم
    \[
    \tan\phi=\frac{0.711538}{\sqrt{1-(0.711538)^2}}
    =\frac{0.711538}{\sqrt{0.4937}}
    =\frac{0.711538}{0.7026}
    =1.0127.
    \]
\end{itemize}
این مقدار همان عددی است که در ستون \(\tan\phi\) جدول \ref{tab:RL_raw} گزارش شده است. از همین روند برای تمام نقاط استفاده شده است.

\paragraph{نکته دربارهٔ گرد کردن}
در محاسبات میانی، اعداد با دقت بیشتری نگه داشته شدند و تنها در ارائهٔ جدول‌ها گرد شدند تا خطای گرد کردن روی شیب و عدم‌قطعیت کمینه شود. در گزارش نهایی، مقدارها با تعداد رقم معنی‌دار سازگار با عدم‌قطعیت گزارش شدند (برای مثال \(L=(0.69\pm 0.06)\,\mathrm{H}\)).



\section{ضمیمهٔ H: نکات اجراییِ اسیلوسکوپ در حالت \lr{XY}}

\begin{itemize}
    \item در حالت \lr{XY}، \lr{Trigger} عملاً نقش پایداریِ زمان‌دامنه را ندارد؛ پایدار بودن شکل به ثابت بودن فرکانس‌ها و دامنه‌ها وابسته است.
    \item برای جلوگیری از اعوجاج، باید از اشباع ورودی‌ها پرهیز کرد و ورودی‌ها را روی \lr{AC Coupling} یا \lr{DC Coupling} مطابق شرایط مدار تنظیم نمود.
    \item بهتر است ابتدا دامنهٔ هر کانال جداگانه در حالت زمان‌دامنه تنظیم شود تا \(A\) و \(B\) دقیقاً قابل خواندن باشند، سپس وارد حالت \lr{XY} شویم.
\end{itemize}


\end{document}
